% =====================================================================
%  Plantilla LaTeX para trabajos de Ingeniería - Universidad Carlos III de Madrid
%  Compilar con: pdflatex -> biber -> pdflatex -> pdflatex  (o Overleaf)
% =====================================================================
\documentclass[12pt,a4paper,twoside,openright]{report}

% --------------------- Codificación e idioma ---------------------
\usepackage[utf8]{inputenc}
\usepackage[T1]{fontenc}
\usepackage[spanish,es-tabla,es-nodecimaldot]{babel}
\usepackage{lmodern}
\usepackage{microtype}

% --------------------- Márgenes y espaciado ----------------------
\usepackage[a4paper,top=3cm,bottom=3cm,inner=3cm,outer=2.5cm]{geometry}
\usepackage{setspace}
\onehalfspacing
\setlength{\parskip}{0.5em}

% --------------------- Colores corporativos ----------------------
\usepackage{xcolor}
\definecolor{azulUC3M}{RGB}{0,0,102}

% --------------------- Matemáticas -------------------------------
\usepackage{amsmath,amssymb,amsthm}
\usepackage{siunitx}
\sisetup{output-decimal-marker={,}}

% --------------------- Figuras y tablas --------------------------
\usepackage{graphicx}
\usepackage{float}
\usepackage{caption}
\usepackage{subcaption}
\usepackage{booktabs}
\usepackage{multirow}
\usepackage{tabularx}
\captionsetup{font=small,labelfont={bf,color=azulUC3M}}

% --------------------- Código fuente -----------------------------
\usepackage{listings}
\lstset{
  basicstyle=\ttfamily\small,
  keywordstyle=\color{azulUC3M}\bfseries,
  commentstyle=\color{gray}\itshape,
  stringstyle=\color{teal},
  numbers=left, numberstyle=\tiny\color{gray},
  frame=single, rulecolor=\color{gray!50},
  breaklines=true, showstringspaces=false, tabsize=4,
  literate={á}{{\'a}}1 {é}{{\'e}}1 {í}{{\'i}}1 {ó}{{\'o}}1 {ú}{{\'u}}1 {ñ}{{\~n}}1
}
\renewcommand{\lstlistingname}{Código}
\renewcommand{\lstlistlistingname}{Índice de códigos}

% --------------------- Cabeceras y pies --------------------------
\usepackage{fancyhdr}
\pagestyle{fancy}
\fancyhf{}
\fancyhead[LE]{\small\leftmark}
\fancyhead[RO]{\small\rightmark}
\fancyfoot[C]{\thepage}
\renewcommand{\headrulewidth}{0.4pt}
\fancypagestyle{plain}{\fancyhf{}\fancyfoot[C]{\thepage}\renewcommand{\headrulewidth}{0pt}}

% --------------------- Títulos de capítulo -----------------------
\usepackage{titlesec}
\titleformat{\chapter}[display]
  {\normalfont\huge\bfseries\color{azulUC3M}}
  {\chaptertitlename\ \thechapter}{15pt}{\Huge}
\titlespacing*{\chapter}{0pt}{-20pt}{30pt}

% --------------------- Bibliografía (estilo IEEE) ----------------
\usepackage{csquotes}
\usepackage[backend=biber,style=ieee]{biblatex}
\addbibresource{referencias.bib}

% --------------------- Otros -------------------------------------
\usepackage{enumitem}
\usepackage[acronym,toc,nomain]{glossaries}
\makeglossaries
\usepackage[hidelinks]{hyperref}   % siempre de los últimos

% --------------------- Acrónimos ---------------------------------
\newacronym{uc3m}{UC3M}{Universidad Carlos III de Madrid}
\newacronym{api}{API}{Application Programming Interface}

% =====================================================================
%  DATOS DEL TRABAJO  (edita aquí)
% =====================================================================
\newcommand{\titulo}{Vive tu municipio}
% Un autor por línea, separados por \\
\newcommand{\autores}{%
  Paula Pato Gómez \\
  Diego Gómez Terradillos \\
  Daniel Muñoz Martínez \\
  Anna de Vismes \\
  Jorge Rodríguez Moreno}
\newcommand{\asignatura}{Datos Masivos y Encadenados}
\newcommand{\titulacion}{Máster en Ingeniería Informática}
\newcommand{\escuela}{Escuela Politécnica Superior}
\newcommand{\lugarfecha}{Leganés, \today}
\newcommand{\curso}{Curso 2026--2027}

\begin{document}

% =====================================================================
%  PORTADA
% =====================================================================
\begin{titlepage}
  \centering
  % Logo oficial UC3M (logo_uc3m.pdf, convertido desde Logo_UC3M.svg)
  \IfFileExists{logo_uc3m.pdf}{\includegraphics[width=4cm]{logo_uc3m.pdf}}{%
    \IfFileExists{logo_uc3m.png}{\includegraphics[width=4cm]{logo_uc3m.png}}{%
      {\Large\bfseries\color{azulUC3M} Universidad Carlos III de Madrid}}}\par
  \vspace{1cm}
  {\large \escuela}\par
  {\large \titulacion}\par
  \vspace{2.5cm}
  {\color{azulUC3M}\rule{\textwidth}{1.5pt}}\par\vspace{0.5cm}
  {\Huge\bfseries\color{azulUC3M} \titulo}\par\vspace{0.5cm}
  {\color{azulUC3M}\rule{\textwidth}{1.5pt}}\par
  \vspace{2cm}
  {\large \asignatura}\par
  \vfill
  \begin{tabular}{rl}
    \textbf{Autores:} & \begin{tabular}[t]{@{}l@{}}\autores\end{tabular} \\[0.2cm]
    \textbf{Curso:}   & \curso
  \end{tabular}\par
  \vspace{1.5cm}
  {\large \lugarfecha}\par
\end{titlepage}

% Página de licencia (habitual en la UC3M)
\thispagestyle{empty}
\vspace*{\fill}
\noindent Esta obra se encuentra sujeta a la licencia Creative Commons
\textbf{Reconocimiento -- No Comercial -- Sin Obra Derivada}.
\cleardoublepage

% =====================================================================
%  CUERPO
% =====================================================================
\chapter{Introducción}
\label{cap:introduccion}

\section{Motivación}
Contexto y problema que se aborda. Ejemplo de acrónimo: \gls{uc3m}.
Ejemplo de cita bibliográfica \cite{knuth1984}.

\section{Objetivos}
\begin{enumerate}[label=\textbf{O\arabic*.}]
  \item Objetivo general.
  \item Objetivo específico.
\end{enumerate}

\section{Estructura del documento}
En el Capítulo~\ref{cap:estado} se revisa el estado del arte; en el
Capítulo~\ref{cap:diseno}, el diseño de la solución, etc.

\chapter{Estado del arte}
\label{cap:estado}
Revisión de trabajos y tecnologías relacionadas \cite{lamport1994}.

\chapter{Diseño e implementación}
\label{cap:diseno}

\section{Arquitectura}
\begin{figure}[H]
  \centering
  \fbox{\parbox[c][4cm][c]{0.7\textwidth}{\centering Aquí va tu figura\\
    \texttt{\textbackslash includegraphics[width=0.7\textbackslash textwidth]\{img/arquitectura\}}}}
  \caption{Arquitectura general del sistema.}
  \label{fig:arquitectura}
\end{figure}


% =====================================================================
%  BIBLIOGRAFÍA Y ANEXOS
% =====================================================================
\cleardoublepage
\phantomsection
\addcontentsline{toc}{chapter}{Bibliografía}
\printbibliography[title={Bibliografía}]

\appendix
\chapter{Anexo de ejemplo}
\label{anx:ejemplo}
Material complementario.

\end{document}
